% Transcription — fonds Alexandre Grothendieck, shelfmark 127, batch 1 (pages 1-20).
% Established from the facsimile archives/batches/127/batch-01.pdf, prepared
% from a copy of 127.pdf fetched through the project's relay (PDF page k =
% archivists' page k, no cover sheet), per the skill
% transcribe-grothendieck. The folder is twenty-three pages; this is the
% first of two batches (pages 21-23 are batch 2, transcribed by another
% pass; its file was read for notation).
%
% Pass: Opus 5.5 (claude-opus-5-5), 2026-09-26 — first pass, unchecked against the pages by a human.
%
% Dating: the folder is « s.d. » in src/content/catalogue.ts; its inventory
% group « Descentes » (cotes 126 to 133) is dated [avant 1970], and \dating
% says both, per the skill's group rule, exactly as folders 126, 128 and
% 127 batch 2 do. Nothing on these leaves dates them; page 8 is a leaf of
% the EGA IV typescript (§ 21.10), which is consistent with the group's
% date but is not recorded as dating.
%
% Pages skipped: none. Page 1 is a tan cover inscribed in ink, top right,
% underlined: « Application de la descente non plate : platitude d'une
% adhérence schématique, et lemme de Hironaka ». It takes no \page marker
% and its words become the section heading, with a \note.
%
% Pages summarised: 8 and 16.
% - Page 8 is a typed leaf, IV-1131 (?), of the EGA IV typescript: the end
%   of a proof, then n° 21.10 « Factorialité des anneaux locaux réguliers »,
%   Th. 21.10.1 (Auslander-Buchsbaum). The text is published, so it gets one
%   French \note; his ink interventions are given in full there (« Créditer
%   Kaplansky de la démonstration ici », « (Auslander-Buchsbaum) »,
%   « (5.11.6) » -> « I 9.4.5 », a struck clause, the label (21.10.1.1),
%   small letter repairs). It is kept, not skipped, because it carries his
%   hand; it is not a draft of this run.
% - Page 16 is in another hand, not his and not identifiable from the leaf
%   (it is not marked as Murre's or Levelt's): a fair French statement
%   « Lemme 1.4' » restating Th. 1.4. Following the brief's rule for other
%   people's text it gets one French \note summarising it; no intervention
%   of his is certain (the X, Y of the conclusion are gone over in heavier
%   ink, which is said).
% No notes by Murre or Levelt occur in this batch.
%
% Pages transcribed from his hand: 2-7, 9-15, 17-20. Blue ink, fast
% drafting, with pencil additions; pp. 10, 12, 14, 15, 19 and much of the
% margins are in pencil. The formulas and numbered statements carry; much
% of the connective prose, and several heavily struck blocks (pp. 3, 4, 7,
% 9), do not and are left \ill{} — following folder 29 batches 8-9.
% His pagination: 1-7 on pp. 2-7 and 9 (p. 9's « 7 » overwrites an 8);
% pp. 10-20 carry no number of his that is sure. Recorded once per page
% in a \note (« sa page n »).
%
% What the batch is about. An application of non-flat descent (descent of
% flat quotients of a given sheaf) to the flatness of a schematic closure.
% Pages 2-6 (« Platitude d'une adhérence schématique »): 1.1 the problem —
% f : X -> Y, U open (retrocompact) in X, F quasi-coherent on X, G_U a
% quotient of F|U flat over Y; find a quotient G of F extending G_U, flat
% over Y, with G -> i_*(G_U) universally injective; the solution is unique,
% G = Im(F -> i_*(G_U)). Lemma 1.2 (the question is local for fpqc; SGA 1
% VIII), Lemma 1.3 (Y artinian, descent along an epimorphism Y' -> Y),
% Theorem 1.4 (Y, X locally noetherian, local rings with geometrically
% normal formal fibres; conditions (a) on specialisations of associated
% points, (b) solutions after base change to a family Y_i -> Spec Ô_{Y,y}
% with zero intersection of kernels, e.g. the Spec(Ô/m^{i+1}), (c) finite-
% ness; then a coherent flat solution exists; NB: (a), (b) necessary by
% EGA IV 12.1.1.5), its proof by localisation, henselisation and
% completion (p. 4), Cor. 1.5 (equidimensional (S_1) fibres), Cor. 1.6
% (« critère valuatif »: traits Y' finite over Spec O_{Y,y}), Cor. 1.7.
% Pages 7-12: Th. 1.8, on applying Hironaka's lemma — for f locally of
% finite presentation with a)-e) at x (Y reduced, f universally open at the
% maximal generisations, equal fibre dimensions, X_y reduced there and
% (X_y)_red geometrically normal), some open U with U_red -> Y normal; a
% cancelled corollary (2) with conditions (i)-(iiii bis); its proof (p. 9,
% reduction to Y essentially of finite type over Z, constructibility,
% EGA IV 15.2.3, « c'est le cas de Hironaka, OK »), a fair restatement of
% 1.8 with the proof for O_{Y,y} excellent (pp. 10, 12: the open V where f
% is flat with separable fibres, EGA IV 12, 15.2.3, then the valuative
% criterion for flatness of a schematic closure). Page 11: an aside on the
% category of finite k-schemes with the Zariski topology, quasi-coherent
% sheaves M ⊗_Λ Λ' and the failure of Ker u to be one (« manque de
% platitude »). Pages 13, 17, 20: formal sheaves on a site C (categories
% S_0, 𝔖, 𝔐, 𝒩), H^*(C_{/δ}, F) ≃ H^*(Ẑ, F|Ẑ), the lim^(1) exact sequence
% for H^q(X^(i), F) and, p. 20, H^n(S_∞, F) against lim H^n(S_i, F) with
% ML and the lim^(p) spectral sequence. Pages 14, 15, 18, 19: pencil
% notes and sketches for 1.4-1.5 (Remarques 1.4.2, 1.4.3, « Vrai théorème »
% / « Énoncé final »: a solution exists iff it does after every base change
% to an artinian scheme or a trait; a counterexample with X -> Y
% birational; EGA IV 11.12 and 12.2).
%
% Notation fixed once for the batch, consistent with batch 2 and folders
% 126/128: \mathcal{C} for his round C (the site, pp. 11, 13, 17);
% X_{\mathrm{red}}, (X_y)_{\mathrm{red}}; X \times_Y X; \mathrm{Ass};
% i_* for the direct image along U -> X; \widehat{\mathcal{O}}_{Y,y},
% \mathfrak{m}_y; \mathfrak{S}, \mathfrak{M} for his gothic-looking S, M
% (p. 11) and \mathcal{M}, \mathcal{N} where the same letters recur (pp.
% 13, 17); \varprojlim^{(1)}; « Pb » for his underlined abbreviation of
% « problème », set \emph{Pb}. Words he underlines are \emph. His circled
% a), b), c) are written (a), (b), (c). Double and triple arrows in the
% sketch of p. 19 are drawn single, with a \note (no parallel-arrow style
% in the renderer), as batch 2 does. No Amitsur complex or non-abelian H^1
% occurs here; the conductor does not appear.
%
% Readings to check first: p. 2 the long struck margin note; p. 3 the ends
% of (b) and (c) and the struck block of the Dém; p. 4 almost all the
% connective prose (the reduction to X' = Spec B̂); p. 7 the struck
% conclusion of 1.8 and the whole « Cor (2) »; p. 9 the second paragraph
% and the boxed struck bottom, including two lines written upside down;
% p. 11 « revêtements locaux » and « & connexes »; p. 13 most of the
% second paragraph; p. 17 the exponent q versus q-1; p. 8 the typed page
% number. Apparatus: 242 \ill{}, 98 \uncertain{}.
% Batch 2 was not contradicted by anything on these pages.

\documentclass[11pt,a4paper]{article}
% Le préambule commun à toutes les transcriptions du fonds Grothendieck.
%
% Il tient en une idée : **l'incertitude se compose, elle ne se résout pas.**
% Une transcription qui lisse un mot douteux a détruit la seule chose pour
% laquelle on la faisait — un lecteur doit pouvoir savoir, à chaque endroit,
% ce qui était sur le feuillet et ce qui a été ajouté. Les six macros qui
% suivent sont donc l'appareil critique, et non de la décoration.
%
% Les mêmes commandes sont reconnues par `scripts/render.mjs`, qui produit la
% vue de lecture du site. Une macro ajoutée ici doit l'être aussi là-bas,
% faute de quoi le rendu échouera bruyamment — ce qui est voulu : un rendu
% silencieusement faux est pire qu'un rendu absent.

\ProvidesPackage{grothendieck}[2026/08/08 Transcriptions du fonds Grothendieck]

\usepackage{amsmath,amssymb,amsthm}
\usepackage{mathrsfs}
\usepackage{stmaryrd}
% Les diagrammes commutatifs sont partout dans ce fonds : c'est la langue
% même de la géométrie algébrique de Grothendieck, et une transcription qui
% les rendrait en prose ne transcrirait rien.
\usepackage{tikz-cd}
% Français par défaut — c'est la langue des feuillets — mais l'anglais est
% chargé aussi : la traduction et le résumé sont des documents anglais, et la
% typographie française (espaces avant les deux-points) y serait une faute.
% Ces documents basculent par \selectlanguage{english} en tête de corps.
\usepackage[english,french]{babel}
\usepackage{fontspec}
\usepackage{geometry}
\geometry{a4paper,margin=25mm,marginparwidth=28mm}
\usepackage{marginnote}
\usepackage{xcolor}
% ulem, not soul. `soul`'s \st and \ul are built on a character-by-character
% scan that XeLaTeX's font machinery breaks: the document fails with an
% undefined control sequence rather than degrading. ulem does the same two jobs
% and is Unicode-engine safe. `normalem` stops it hijacking \emph, which the
% transcriptions use for Grothendieck's own emphasis.
\usepackage[normalem]{ulem}
\usepackage{hyperref}
\hypersetup{colorlinks=true,linkcolor=black,urlcolor=black,pdfborder={0 0 0}}

% --- Métadonnées du lot -----------------------------------------------------
% Reprises telles quelles par le script de rendu, qui les affiche en tête de
% la vue de lecture. Elles fixent ce que le fichier prétend couvrir : sans
% elles, une transcription flotte sans point d'attache dans un fonds de seize
% mille feuillets.

\newcommand{\folder}[1]{\gdef\grothFolder{#1}}
\newcommand{\batch}[1]{\gdef\grothBatch{#1}}
\newcommand{\leaves}[2]{\gdef\grothFirst{#1}\gdef\grothLast{#2}}
\newcommand{\foldertitle}[1]{\gdef\grothFoldertitle{#1}}
\gdef\grothFolder{?}\gdef\grothBatch{?}\gdef\grothFirst{?}\gdef\grothLast{?}\gdef\grothFoldertitle{}

% --- Le résumé --------------------------------------------------------------
%
% Il ouvre la lecture modernisée, sous le titre « Résumé », et il est composé
% en petit. Ce n'est pas un ornement : c'est le seul passage du document qui
% ne soit pas des mathématiques, et un lecteur doit pouvoir le reconnaître
% avant de l'avoir lu — puis le sauter s'il connaît déjà le sujet, ou s'y
% arrêter s'il ne le connaît pas. Le filet qui le suit marque la reprise du
% corps.

\newenvironment{resume}
  {\small\setlength{\parskip}{0.45em}}
  {\par\medskip\hrule\medskip}

% --- Mots-clés ---------------------------------------------------------------
%
% En anglais, en clôture du résumé de la lecture modernisée : le vocabulaire
% moderne sous lequel chercher ce que les feuillets construisent. Le manifeste
% du site (`npm run manifest`) les extrait pour étiqueter la cote entière —
% cette ligne est la source unique des tags, il n'y a pas de fichier à côté.
\newcommand{\keywords}[1]{\par\smallskip\noindent
  {\footnotesize\textsc{Keywords} --- \textit{#1}}\par}

% --- L'appareil critique ----------------------------------------------------

% Le numéro de feuillet, en marge. C'est l'ancre qui permet de régler un
% désaccord sur une formule en regardant *un* feuillet plutôt qu'une chemise
% de six cents. Le site s'en sert aussi pour tourner les pages du fac-similé
% quand on fait défiler la transcription.
\newcommand{\leaf}[1]{%
  \marginnote{\footnotesize\color{gray!70}\textsf{#1}}%
  \label{leaf:#1}\ignorespaces}

% Illisible : on ne devine pas. Un mot inventé qui se lit comme les autres est
% le pire résultat possible.
\newcommand{\ill}{\textcolor{red!60!black}{[\,\ldots\,]}}

% Lecture douteuse : le mot est donné, et signalé comme incertain.
\newcommand{\uncertain}[1]{\uline{#1}}

% Ajout de l'éditeur — une lettre manquante, un mot sous-entendu.
\newcommand{\add}[1]{\textcolor{blue!50!black}{[#1]}}

% Biffé par Grothendieck lui-même. Ce qu'il a rayé fait partie du document :
% une rature dit souvent où la pensée a changé de direction.
\newcommand{\struck}[1]{\sout{#1}}

% Note du transcripteur, sur l'état matériel du feuillet.
\newcommand{\note}[1]{\par\smallskip{\small\color{gray!60!black}\textit{[#1]}}\par\smallskip}

% Note marginale de Grothendieck, distincte du corps du texte.
\newcommand{\marginal}[1]{\par\smallskip\noindent\hspace{1em}%
  {\small\color{gray!60!black}#1}\par\smallskip}

% --- Titre ------------------------------------------------------------------

\AtBeginDocument{%
  \begin{center}
    {\large\bfseries Fonds Alexandre Grothendieck --- cote n\textsuperscript{o}~\grothFolder}\\[2pt]
    {\itshape\grothFoldertitle}\\[4pt]
    {\small feuillets \grothFirst--\grothLast\ (lot \grothBatch)}\\[2pt]
    {\footnotesize\color{gray!60!black}Transcription établie d'après le fac-similé numérisé
     par l'Université de Montpellier.\\
     Les lectures incertaines sont soulignées, les illisibles notées [\,\ldots\,],
     les ajouts entre crochets.}
  \end{center}
  \vspace{1em}}

\endinput


\folder{127}
\batch{1}
\pages{1}{20}
\dating{s.d. — le groupe « Descentes » (126 à 133) est daté [avant 1970]}
\watermark{Édition de démonstration}
\foldertitle{Descente non plate : notes manuscrites (s.d.)}

\begin{document}

\section*{Application de la descente non plate : platitude d'une adhérence schématique, et lemme de Hironaka}
\note{inscrit à l'encre, souligné, en haut à droite du feuillet de garde
(page 1), qui ne porte rien d'autre ; le mot lu « lemme » est repassé
ou surchargé. Le feuillet ne reçoit pas de numéro de page}

\subsection*{Platitude d'une adhérence schématique}

\page{2}
\note{sa page 1 ; le titre est souligné}
\marginal{$\overline{U \cap T}$ \quad $|X|$}\note{au crayon, en haut à
droite ; la disposition (un trait, $|X|$ dessous) n'est pas claire}

\emph{1.1.} $U \subset X$, $\mathbf{F}$, $f \colon X \to Y$.
Données $f \colon X \to Y$, $U$ ouvert \add{(}\uncertain{rétrocompact}\add{)}
de $X$, $F$ quasi-cohérent sur $X$, $G_U$ quotient de $F|U$,
quasi-cohérent et \emph{plat} sur $Y$.
\note{« $X \to Y$ » et « rétrocompact » (au crayon, encerclé) sont
ajoutés en interligne}

\emph{Pb} Trouver un quotient \add{quasi-coh.}\ $G$ de $F$, qui prolonge
$G_U$, qui soit plat sur $Y$, et tel que
$G \to G_U$ \marginal{(rel.\ à $i$)} soit \emph{univ.\ injectif} rel/$S$
\add{[$G \to i_*(G_U)$]}.
\note{« quasi-coh. » est en interligne ; « [$G \to i_*(G_U)$ » est écrit
au-dessus de « univ.\ injectif », et « rel.\ à $i$ », au crayon et
encerclé, est relié par un trait à $G \to G_U$}
[Il \uncertain{faux et} suffit pour ceci, \struck{\ill{}} si $f$ est
loc.\ de présentation finie, et si \struck{$Z = S_U G$, que} $G$ est de
prés.\ finie sur $S$, que pour tt $y \in Y$, $U_y$ contienne
$\mathrm{Ass}\, G_y$].
\note{les deux mots au-dessus de « suffit », au crayon, sont peu sûrs}

\marginal{\ill{} \ill{} \uncertain{peut supposer} $X$ plat sur $S$
\ill{} sur $S$ si à l'origine il \uncertain{était} \ill{} de type fini)
et $F$ \struck{plat} sur $X$ (et de type fini s'il l'était
initialement), quitte à \ill{} \ill{}}
\note{note écrite en long dans la marge gauche, barrée de traits
obliques ; le début est illisible}

Si le \emph{Pb} a une solution, elle est unique,
\[
  G \simeq \mathrm{Im}\bigl(F \to i_*(G_U)\bigr) .
\]
\note{sous la flèche, un ajout : $\to i_*(F_U)$, avec une flèche
remontant vers $i_*(G_U)$}
La question est donc de savoir si le $G$ défini précédemment \add{est}
plat rel/$Y$ et $G \to G_U$ est \struck{\ill{} \ill{}} univ.\ injectif
rel/$Y$. \add{(}\emph{Remarque} Question \emph{locale} sur $X$ et $Y$,
pour la \uncertain{topologie} \ill{}\,!\add{)}
\uncertain{L'unicité s'explique} \ill{}.
\note{la remarque est ajoutée en interligne, encadrée d'un trait}

\emph{Lemme 1.2.} Question locale \add{sur $Y$} pour fpqc. De même la
question analogue, où on veut que $G$ soit loc.\ de présentation finie sur
$X$. // Plus gén., descente possible sur $X$ par morphismes fid.\ plat
qu.\ cpct $X' \to X$.
\marginal{SGA 1 VIII}\note{encerclé, dans la marge gauche, en face du
lemme 1.2}

La descente non plate \add{((des quotients \emph{plats} d'un faisceau
déjà donné))} permet aussi de prouver :
\note{la double parenthèse est ajoutée au crayon en interligne}

\emph{Lemme 1.3.} Supposons $Y$ artinien
\add{$= \mathrm{Spec}\, A$}, et qu'il existe un épimorphisme de schémas
$Y' \to Y$, tel que après ce changement de base, le \emph{Pb} ait une
solution. Suppo\ldots
\marginal{$A = \Gamma(Y, \mathcal{O}_Y)$}
\marginal{Il faut supposer ici $Y' \to Y$ fini \ill{} et $f \colon X \to
Y$ loc.\ \struck{de prés.}\ \ill{}, ou alors de prés.\ finie}
\note{« $A = \Gamma(Y, \mathcal{O}_Y)$ » est encerclé et relié à
« artinien » ; la seconde note, au crayon, est écrite en oblique dans
le coin inférieur gauche et chevauche le mot « Lemme »}

\page{3}
\note{sa page 2}

$Y' \to Y$ fini, \emph{ou} $X$ loc.\ de prés.\ finie, et la solution
$G'$ de présentation finie. Alors le \emph{Pb} a une solution,
\struck{\ill{}} $Y$, et dans le 2\textsuperscript{ème} cas, la solution $G$
est aussi de présentation finie.

\emph{Théorème} \struck{Lemme} \emph{1.4.} \add{Néc.\ Suff.}
Supposons que $X$, \struck{\ill{}} $Y$ soient loc.\ noeth.,
\struck{\ill{}} \add{\uncertain{(loc.\ intègres)}} les anneaux locaux de $X$
à fibres formelles géom.\ \struck{\ill{}} \ill{} \add{($F$, $G$
cohérents)}.
\note{« Théorème 1.4 » est encadré, « Lemme » biffé dessous ; « Néc.\
Suff. » est ajouté au crayon et relié à l'énoncé ; « $X$, » est en
interligne au-dessus d'un mot biffé}
\marginal{en tout pt de $X - U$}\note{encadré au crayon, dans la marge
gauche}

Supposons que (a) pour tout couple de points $y \neq y'$ de $Y$, avec
$y \in \overline{\{y'\}}$, \add{et tt générisation} $x \in X_y - U_y$
\struck{$x \in U_y$ tel q} \add{($x' \in \mathrm{Ass}(G_U)_{y'}$ de $x$)},
il existe une spécialisation $x_1$ de $x'$ \struck{qui} qui est dans
$U_y$, et qui généralise $x$, --- et que la condition \uncertain{soit
vérifiée} pour tt $x'$ \ill{} $X$.
\note{a), b), c) sont cerclés ; les ajouts sont en interligne, serrés et
en partie superposés}

\begin{tikzcd}[column sep=small, row sep=small]
U_y \ni x_1 \arrow[d] & \\
X_y \ni x & x' \arrow[ul] \arrow[l] \\
y & y' \arrow[l]
\end{tikzcd}

\note{croquis de la marge gauche : $x'$ se spécialise en $x_1$ et en $x$,
$x_1$ en $x$, et $y'$ en $y$ ; à côté, au crayon :
$T = \overline{\{x'\}}$, $(T_y \cap U) \cup (T_y \cap Z) \subset
\uncertain{T_y}$}

(b) Supposons de plus que pour tt $y \in Y$, il existe une famille de
morphismes locaux
\[
  Y_i \to Y' = \mathrm{Spec}(\widehat{\mathcal{O}}_{Y,y}),
\]
\add{les $Y_i$ noeth.\ locaux,} tels que après changement de base le
\emph{Pb} ait une solution \struck{\ill{} \ill{} \ill{}}
\add{$\bigcap_i \mathrm{Ker}(\widehat{\mathcal{O}}_{Y,y} \to
\mathcal{O}_{Y_i,y_i}) = 0$}. (c) \uncertain{Supposons} de plus,
\ill{} $f$ loc.\ de t.\ finie, \struck{\ill{}} \add{\uncertain{(les)}}
$Y_i \to Y'$ finis. Alors le \emph{Pb} a une solution, qui est un Module
cohérent (plat sur $Y$).
\note{la fin de (b) et le début de (c) sont très surchargés ; « cohérente »
est ajouté puis biffé au-dessus de la ligne}

\emph{NB} Si $X$ loc.\ de type fini $/S$, \struck{\ill{}} les cond.\ a) et
b) \uncertain{sont} \emph{nécessaires} pour l'existence d'une solution,
en vertu de IV 12.1.1.5. Dans le cas \uncertain{où} $X$ \uncertain{est}
\ill{} \ill{} les anneaux locaux \ill{} \ill{}
\add{\uncertain{la condition gén.\ \ill{} pour \ill{} de $Y$}}.
\note{la NB est écrite serrée, en interligne à droite ; « IV 12.1.1.5 »
est net}
\uncertain{Utiliser} pour b), et pour c) \uncertain{utiliser} que les
conditions \uncertain{sont} \ill{} \ill{} 1.3.

On peut prendre pour $Y_i$ les
\[
  \mathrm{Spec}\bigl(\widehat{\mathcal{O}}_{Y,y} / \mathfrak{m}_y^{i+1}\bigr),
  \qquad i \in \mathbb{N} .
\]

\emph{Dém} \struck{On se localise en \ill{} l'image dans $X$ des}
On se localisant en un point maximal de \add{l'ens.\ des points
$x \in X_y - U_y$ \struck{\ill{}} ($y = f(x)$) qui est dans
$\mathrm{Ass}\, G_y$} \struck{\ill{}} qui admette une gén.\ dans
$X_y - U_y$, on peut supposer que $X \to Y$ est local,
\add{$G = \mathrm{Im}\bigl(F \to i_*(G_U)\bigr)$}
\struck{\ill{} plat, \ill{} peut supposer \ill{} et que $G$ est plat
$/Y$ \ill{} \ill{} $Y_i$ \ill{} fermé de $X$. De plus, par descente plate,
\ill{} $x$ de $X_y$ \ill{} \ill{} \ill{} \ill{} de platitude, \ill{} \ill{}
$X$ est local et que $x \in X_y$ \ill{}}.
On peut supposer $U = X - \{x\}$.
\struck{[NB. $\mathrm{Ass}\, G_Y = \mathrm{Ass}\, G_U$] $= \mathrm{Ass}\, G$}
\note{le bloc du bas est encadré, barré ligne à ligne et traversé au
crayon d'une ligne ondulée ; les ajouts au-dessus sont en interligne}

\marginal{\emph{Remarque} : \ill{} $\exists$ solution si \ill{}
\uncertain{trivial} \ill{} \ill{} dans le complémentaire du plus grand ouvert
au-dessus duquel $G$ \ill{} les propriétés \ill{} \ill{} on peut supposer
$X \to Y$ local, $U = X - \{x\}$ ($x$ point fermé de $X$)}
\marginal{\ill{} $\forall x' \in \mathrm{Ass}\, G_{y'}$, pour tout
$y \in Y$, \ill{} dans $T_y$}
\note{deux notes au crayon dans la marge gauche, la première en oblique,
la seconde en long et encadrée ; « trivial » est encerclé}

\page{4}
\note{sa page 3}

\add{Par descente plate,} on peut supposer $X$ local
\struck{\uncertain{complet}} \add{hensélien}. \struck{\ill{}}
\uncertain{Mais} si \ill{} \ill{} \ill{} \struck{\ill{} solutions $G$}
\ill{} l'hypothèse \ill{} \ill{} Ass \ill{} \uncertain{conservées}
\uncertain{par} (a) (b) \ill{}
\note{le haut de la page est un passage raturé, en partie encadré ; les
mots lus ne sont donnés qu'à titre indicatif}
\ldots\ il \uncertain{est évident} que les hypothèses faites au départ,
\struck{\ill{}} \struck{\uncertain{(puisque} $X^h = \mathrm{Spec}\, B^h$)
\ill{} \ill{}}
l'hypothèse que $B = \mathcal{O}_{X,x}$ est à fibres formelles
géom.\ normales (ce qui est en effet une condition stable par
hensélisation). Par contre, \ill{} \ill{} que le morphisme
$X' = \mathrm{Spec}\, \widehat{B} \to X = \mathrm{Spec}\, B$ est à
fibres \struck{formelles} \add{(géom.)} intègres, donc si $Z_\alpha$
sont les cycles associés à $G$ sur $X$, \struck{et $X'$} dans cette
\ill{}, les $Z'_\alpha$ sont les cycles associés à $G'$ sur $X'$ (qui
sont réunion \ill{} \add{$Z'_\alpha \cap U \neq \emptyset$ pour tout
$\alpha$ ; \uncertain{sont alors vérifiée}} anneaux intègres\,!). Donc
l'hyp.\ (a) \ill{} \ill{} $(X, G, U)$ en $(X', G', U')$. Cela nous
\uncertain{ramène} en remplaçant $X$ par $X'$, l'hyp.\ (b) \uncertain{restant}
\uncertain{vérifiée} (\ill{} \ill{} \uncertain{des solutions} \ill{}
vérifiée). Utilisant la \uncertain{suppl.}\ des \ill{} \ill{} $Y_n$, et
prenant \ill{} la \ill{} \emph{Pb} relatif pour les \ill{} \ill{} des
\uncertain{modules} quotients \uncertain{limite}, on trouve un Module
quotient $G \to H$ de $G$, qui est plat sur $Y$, et qui \ill{}
\ill{} sur les \ill{} tel que l'homom.\ $G \to H$ soit \ill{} \ill{}
pour tt $Z_\alpha$, \ill{} $Z_\alpha \cap U_y \neq \emptyset$, \ill{}
pts de $U_y$ ; comme pour $G \to H$ est injectif sur $U$, donc on conclut
que $G \to i_*(G_U)$ est injectif. Donc injectif puisque $G$ est
\emph{plat} sur $Y$, on a $G = H$, i.e.\ $G$ est plat sur $Y$, et tel
que $x \notin \mathrm{Ass}\, G_y$, donc vu que $G_y = H_y$ \ill{}
\ill{} \ill{}.
\note{« il est évident » et « Par contre » ouvrent des lignes serrées dont
la liaison n'est pas sûre ; un triple trait vertical en marge gauche
marque l'alinéa « \ldots\ intègres, donc si $Z_\alpha$ \ldots »}

\begin{tikzcd}[column sep=small, row sep=small]
x_1 \arrow[d] & x' \arrow[l] \arrow[dl] & \hat{x}_1 & \hat{x}' \in \mathrm{Ass}\, \widehat{U} \\
x & & & \hat{x} \\
X_n \arrow[r] & X \arrow[d] & \widehat{X} = X' \arrow[l] & \supset U' \\
Y_n \arrow[r, hook] & Y & & \\
y & y' \arrow[l, no head] & &
\end{tikzcd}

\note{croquis de la marge gauche, entouré d'une accolade avec $G_n$,
$U_n$ ; les flèches entre $\hat{x}_1$, $\hat{x}'$ et $\hat{x}$ sont
surchargées et ne sont pas reproduites ; un mot biffé suit
« $\mathrm{Ass}\, \widehat{U}$ »}

\begin{tikzcd}[column sep=small, row sep=small]
F_n \arrow[r, "u_n"] \arrow[d, "\text{surj}"'] & i_{n*}(G_{U_n}) \\
G_n \arrow[ur] \arrow[r] & H_n
\end{tikzcd}

\note{au crayon dans la marge gauche, entouré ; dessous :
« $H_n = \mathrm{Im}\, u_n = \mathrm{Im}\, v_n$ est un quotient de
$G_n$ » ; l'étiquette $v_n$ n'est pas placée sur une flèche lisible}

\page{5}
\note{sa page 4 ; la moitié supérieure (le Cor.~1.5) est traversée au
crayon d'une grande croix, et une ligne au crayon la sépare du reste}

\emph{Cor.~1.5.} Supposons que $f$ soit loc.\ de type fini,
\add{$F$, $G$ cohérents} et que (a') les $(G_U)_y$ soient \add{$(S_1)$
et} équidimensionnels de dim.\ donnée $d$. \add{[définissant encore $G$
comme $\mathrm{Im}\bigl(F \to i_*(G_U)\bigr)$]} tandis que
\struck{que $U$ soit schém.\ dense} pour tt $y$, $U_y$ soit dense dans
$X_y$. // \struck{Alors} Supposons de plus que pour $x \in X - U$
\struck{\ill{}} l'anneau local $\mathcal{O}_{X,x}$ soit à fibres
formelles géom.\ normales. \struck{Alors la cond} (il suffit que ce
soit le cas de $\mathcal{O}_{Y,f(x)}$\,!). Alors la condition (b) du
th.\ est néc.\ et suffisante pour l'existence d'une solution du
\emph{Pb}.
\marginal{\uncertain{S-M} $G = X$ ?}\note{au crayon, encerclé et barré,
dans la marge gauche ; au-dessous, un gribouillis au crayon}

\emph{Critère valuatif} :

\emph{Cor.~1.6.} Notation et hyp.\ préliminaires de 1.4 \add{(y compris
l'hyp.\ des fibres formelles)}. Supposons de plus que pour tt
$x \in X - U$, posant $y = f(x)$, $\mathcal{O}_{Y,y}$ soit
\struck{\ill{}} réduit, et que pour tt changement de base $Y' \to Y$,
\add{trait} \struck{\ill{}} \add{($=$ spectre d'un anneau de val.\
discrète)} avec $Y'$ fini sur $\mathrm{Spec}\, \mathcal{O}_{Y,y}$, on
ait une solution du \emph{Pb}. Alors on a une solution du \emph{Pb}.
\marginal{1.5\,? \ill{}, mais sans \uncertain{hypothèse} du type de
1.4 a)}
\marginal{Fait double emploi avec 1.7, sauf qu'on a fait
d'hyp.\ \uncertain{« de type fini »}}
\note{deux notes au crayon dans la marge gauche, la première sous un
« 1.5\,? » encerclé ; au-dessus de la seconde, un « $\wedge$\,? » isolé}

\emph{Cor.~1.7.} $f \colon X \to Y$ \struck{\ill{}}
\add{$Y$ loc.\ noeth.}, \struck{\ill{}} $F$, $G_U$ cohérents
\struck{\ill{}}. On suppose :

a) Pour tt $x \in X - U$, posant $y = f(x)$, $\mathcal{O}_{Y,y}$ est
\struck{\ill{}} réduit, et à fibres formelles géom.\ \struck{normales}
\struck{\ill{}} normales \add{(\uncertain{cf.}\ \ill{} de type fini sur
$\mathrm{Spec}\, \mathbb{Z}$)}.

b) Pour $y$ comme dessus, et tt \emph{trait} $Y'$ fini sur
$\mathrm{Spec}(\mathcal{O}_{Y,y})$ \add{\uncertain{et loc.}}, il y a une
solution du \emph{Pb} après chgt de base $Y' \to Y$. \add{Explicitez}

\struck{c) \ill{} de 1.4 \ill{} \ill{} \ill{} conditions \ill{} de 1.5}
\note{la dernière ligne est barrée d'un zigzag au crayon}
\marginal{NB Dans cette condition b) l'hypothèse de platitude devient
triviale, seule la question des $\mathrm{Ass}\, G'_{y'}$ \ill{} est
sérieuse}\note{au crayon, dans la marge gauche, en face de b)}

\page{6}
\note{sa page 5 ; deux lignes seulement, en haut du feuillet}

\emph{Alors} le \emph{Pb} posé admet une solution.

Nous allons maintenant utiliser le \emph{lemme de Hironaka}.

\subsection*{Th.~1.8}

\page{7}
\note{sa page 6}

\emph{Th.~1.8.} $f \colon X \to Y$ loc.\ de prés.\ finie.

(1) Soient $x \in X$, $y = f(x)$.
\note{le chiffre devant « Soient » est cerclé et surchargé ; la lecture
(1) s'appuie sur le « (2) » cerclé plus bas}

a) $Y$ réduit en $y$

b) $f$ univ.\ ouvert en les gén.\ maximales de $x$ dans $X_y$

c) Pour tte gén.\ \struck{\ill{}} $y'$ de $Y$, et toute \add{\uncertain{les}}
générisation $x'$ de $x$ appartenant à $X_{y'}$, les
$\dim_{x'} X_{y'}$ sont égaux à \struck{\ill{}} $\dim_x X_y$

d) $X_y$ \struck{réduit \ill{}} réduit en les gén.\ maximales de $x$
dans $X_y$

e) $(X_y)_{\mathrm{red}}$ est géom.\ \emph{normal} au voisinage de $x$

[d) et e) équivalent à : $\exists$ voisinage ouvert $U$ de $x$ \add{dans
$X_y$} tel que $\overline{U}$ soit géom.\ réduit et
$\overline{U}_{\mathrm{red}}$ normal].
\note{sous $\overline{U}$, un signe « $=$ » vertical et
$U \otimes_{k(y)} \overline{k(y)}$}

\marginal{peut-être remplacer b) par c') $f$ univ.\ ouvert dans un
voisinage de $x$, et c) \uncertain{u.l.} les \uncertain{gén.}\ maximales}
\note{à l'encre, encadré, dans la marge gauche, avec une flèche vers b)}

\marginal{NB Si on suppose $Y$ loc.\ noeth., dans l'énoncé
\uncertain{il revient} à dire que, si $X$ est réduit, alors au
voisinage de $x$ $f$ est normal. Si $X$ et $Y$ sont loc.\ noeth., on
doit pouvoir supprimer l'hyp.\ que $f$ est \struck{loc.\ \ill{}} de prés.\
finie (moyennant des hypothèses d'excellence éventuellement \ill{}
---)}\note{encadré en haut à droite, au crayon puis à l'encre ; les deux
paragraphes sont séparés par un trait}

\emph{Alors} (1) Il existe un voisinage \add{ouvert} $U$ de $x$ tel que
$U_{\mathrm{red}} \to Y$ soit \struck{normal} loc.\ de
\uncertain{présentation} \add{\uncertain{finie et}} \ill{}
\struck{$U \to Y$ \ill{} univ.\ ouvert, à fibres géométriquement
\ill{} \ill{} \uncertain{totalement} \ill{} \ill{} normales, dont
\ill{} \ill{} telles que $(X_y)_{\mathrm{red}}$ normales, \ill{} \ill{}
aussi \ill{} \ill{} \ill{}},
\note{passage barré de traits horizontaux et de trois grands traits
obliques ; lu en partie seulement}

\struck{2) Si $Y$ est réduit, la condition \ill{} aussi \ill{}. Soit
$Y_0 = Y_{\mathrm{red}}$, $X_0 = X \times_Y Y_0$.}

\struck{et si $Y$ est réduit, alors} (i) les conditions \uncertain{aussi
gén.} \ill{}

\emph{Cor} (2) Conditions équivalentes : (i) \struck{\ill{}}
\add{(i bis)} $f$ est univ.\ ouvert, \add{\uncertain{lisse}} et
équidimensionnel \struck{\ill{}} en tt $x \in X$, \struck{\ill{}} ses
fibres sont géom.\ réduites, \struck{\ill{}} et ses fibres
\uncertain{réduites} géom.\ \emph{normales}.

\struck{(ii) $\exists$ un sous-}
(ii) Il existe un ouvert $V$ de $X_0$, tel que pour tt $y \in Y_0$,
$V_y$ soit dense dans $X_y$, et $V \to Y_0$ un morphisme
\struck{$X_0 \to Y_0$ \ill{} à fibres loc.}\ séparable. De plus pour tt
$y \in Y_0$, $(X_{0,y})_{\mathrm{red}}$ est géom.\ normal.
\note{en interligne à droite : « $X_{\mathrm{red}} \to Y_{\mathrm{red}}$ »
et, dessous, « $X_0 \to Y_0$ » biffé, puis « est normal »}

(iii) \struck{Il existe un sous} \ill{} \add{loc.}\ et de présentation
finie

(iiii bis) Il existe un sous-schéma \struck{\ill{}} $Z_0$ de $X_0$,
\add{et loc.\ de prés.\ finie} sur $Y_0$, \struck{et} ayant \ill{} que
\ill{} $X_0$ en les \ill{} \struck{\ill{}} normal, et coïncidant avec
$X_0$ \ill{} \ill{} des fibres \ldots\ \add{relations} \ldots\ normal. De
plus, il existe un ouvert \add{(schématiquement dense)} dans $V$ de $U$
\struck{dans $S$}, tel que $V \to X$ soit \uncertain{univ.}\ ouvert.
\note{toute la partie « Cor (2) » est enfermée dans un cadre de traits
obliques ; « Cor », souligné, est dans la marge gauche, relié au « (2) »
cerclé. Numérotation (i), (i bis), (ii), (iii), (iiii bis) telle qu'écrite}

\page{8}
\note{feuillet dactylographié, qui n'est pas un brouillon de cette
suite : la page IV-11\uncertain{31} (le numéro est complété à la main)
d'un tapuscrit d'EGA IV, § 21 --- la fin d'une démonstration
(un $\mathcal{O}_X$-Module inversible $L$ dont l'image réciproque dans
$X'_0$ est isomorphe à $L_0$, donc ample, d'après (21.9.12), (9.6.4),
(II, 5.3.1 et II, 5.5.3)), puis le n° 21.10 « Factorialité des anneaux
locaux réguliers » : Théorème (21.10.1), un anneau local noethérien
régulier est factoriel, démontré par récurrence sur la dimension $n$ ;
pour $n \geqslant 2$, $X = \mathrm{Spec}(A)$, $U = X - \{a\}$, on est
ramené par (21.6.14) à prouver $\mathrm{Pic}(U) = 0$ ; un
$\mathcal{O}_U$-Module inversible $L$ se prolonge en un
$\mathcal{O}_X$-Module cohérent $F$, qui a une résolution gauche finie
par des $\mathcal{O}_X^{n_i}$, d'où par restriction une résolution
(21.10.1.1) de $L$ ; « le théorème résultera alors des considérations
générales suivantes », et la page s'arrête sur « Sur un espace annelé
$X$, soit $E$ un $\mathcal{O}_X$-Module localement libre ». Texte
publié (EGA IV, 21.10) ; il n'est pas recomposé ici. Interventions à la
main : dans la marge gauche, « Créditer Kaplansky de la démonstration
ici », relié par un trait au début de la démonstration ; en interligne
au-dessus de « Théorème (21.10.1) », « (Auslander-Buchsbaum) », relié
par un crochet ; « factoriel » corrigé (le $l$ final repassé) ;
« (5.11.6) » biffé et remplacé par « I 9.4.5 » ; la proposition
« Puisque $\dim(A) \geqslant 2$ et que $A$ est régulier, » biffée
d'un trait ; le repère « (21.10.1.1) » écrit à la main devant la
seconde résolution ; de petites retouches de lettres (« ramené »,
« résolution », « allons »). Un grand trait vertical au crayon traverse
la page}

\page{9}
\note{sa page 7 (le chiffre surcharge un 8)}

\emph{Dém} \struck{On est ramené par \uncertain{procédé} standard}
\add{La conclusion \uncertain{(ii)} signifie aussi qu'il existe} un
sous-préschéma $U_0$ \struck{\ill{}} de $U$, de prés.\ finie sur $Y_0$,
tel que $U_0 \to U$ soit \emph{surjectif}, \struck{et tel que} tel que
$U_0 \to Y$ \struck{soit géométriquement \ill{} \ill{}} \add{soit
normal} \struck{\ill{} loc.\ de prés.\ finie}. \add{Ceci implique donc}
que $U_0$ est réduit, donc $U_0 = U_{\mathrm{red}}$), et un ouvert $V$
de $U_0$, schématiquement dense rel/$Y$, tel que $V \to X$ soit une
immersion ouverte. Sous cette forme, on voit qu'on est ramené au cas où
$Y$ est le spectre d'un anneau local essentiellement de type fini sur
$\mathbb{Z}$. \struck{\uncertain{À} \ill{} qu'il existe un \ill{}}
\note{la lettre qu'il note $U$ est une capitale surchargée, qui tient du
$U$ et du $V$}

Utilisant la technique \add{de constructibilité et} de passage à une
limite, et le lemme de Hironaka, on voit que l'on est ramené aux
hypothèses sur la fibre $X_{\bar{x}}$ \add{avec $d_0 =$ dim loc.\ d.}
\add{(EGA IV 16.1)} \struck{\ill{}} \ill{} et la \struck{\ill{}}
\add{et équidim.}\ \ill{} \ill{} voisinage. Utilisant la condition
(c), on trouve que $f$ est univ.\ ouvert \add{\uncertain{équidim.}\ en
tout \ill{} de $f$ \ill{} à fibres géom.\ \uncertain{sém.\ réduites}},
en un voisinage de $x$. \struck{D'autre part, \ill{}} et à
\struck{fibres géom.\ \ill{}} \add{telles} que $(X_y)_{\mathrm{red}}$
soit normal. \struck{\ill{}} On \add{peut supposer} $X$ réduit :
\uncertain{prouvons} que $X \to Y$ est normal.
\marginal{utiliser (b) (c) (d) (e)}\note{au crayon, les lettres
cerclées}
\marginal{Lemme à dégager}\note{souligné deux fois, dans la marge gauche,
avec un trait vertical}

Soit $V$ l'ensemble des pts $\bar{x}$ de $X$ tels que $X_y$
\add{($y = f(x)$)} soit \ill{} (a) et réduit (donc séparable) en $x$
\ill{} hyp. En vertu de (a) et \add{EGA IV 15.2.3}, c'est aussi le plus
grand ouvert sur lequel $f$ induit un morphisme \struck{réduit}
séparable, et en vertu de l'hypothèse faite (d, e) $V$ est dense dans
$X$ fibre par fibre.
\struck{Si donc \ill{} $X$ \ill{} \ill{} $X_{\mathrm{red}}$ \ill{}}

\struck{Si $X$ est l'adhérence schématique \add{rel/$Y$} de $V$.
\ill{} \ill{} que $X$ est réduit \ill{} (\ill{} \ill{}), \ill{} on
\uncertain{va} prouver par \ill{} \ill{}}
\ldots\ (par morceaux\,!) \ldots\ c'est le cas de Hironaka, OK.
\note{le bas de la page est encadré et barré de traits obliques ; deux
lignes y sont écrites tête-bêche (la page retournée) et ne sont pas
lues ; « (par morceaux\,!) » et « c'est le cas de Hironaka, OK »
restent hors du cadre}
\marginal{On peut procéder par récurrence sur $\dim Y$, et supposer
\ill{} \uncertain{vrai pour les plus petits}}\note{à l'encre, barré en
croix, dans la marge gauche}
\marginal{\uncertain{(1)} \uncertain{à démontrer}}\note{au crayon, encerclé}
\marginal{Le critère valuatif nous ramène au cas où $Y$ est le spectre
d'un anneau de val.\ discrète}\note{au crayon, en bas à gauche}

\page{10}
\note{sans numéro de sa main ; une reprise de l'énoncé du Th.~1.8, mise
au net}

\begin{tikzcd}[row sep=small]
X \arrow[d, "f\ \text{de t.f.}"] \\
Y
\end{tikzcd}

$x$ au-dessus de $y$, $Y$ \emph{noeth.}
\[
  \left\lbrace
  \begin{array}{l}
    \mathcal{O}_{Y,y} \text{ réduit} \\
    \mathcal{O}_{X,x} \text{ réduit} \\
    f \text{ universellement ouvert (EGA IV 14, 15)} \\
    X_y \text{ réduit en les gén.\ maximales } x_i \text{ de } x \text{ dans } X_y \\
    (X_y)_{\mathrm{red}} \text{ est géom.\ normal}/k(y)
  \end{array}
  \right.
\]
\note{les trois premières lignes et les deux dernières sont réunies par
deux accolades distinctes ; en interligne, sous « $\mathcal{O}_{X,x}$
réduit » : « $\exists\, d$ tel que, au voisinage de $x$, tous les comp.\
\uncertain{irr.}\ des fibres de $f$ soient de dim.\ $d$ »}

les deux dernières conditions $\Longleftrightarrow$
$\overline{X_y}$ $\bigl(\overset{\text{df}}{=} X_y \otimes_{k(y)}
\overline{k(y)}\bigr)$ est, au voisinage de $\bar{x}$, géom.\ réduit,
et $(\overline{X_y})_{\mathrm{red}}$ est normal
\[
  \Downarrow
\]
$\exists$ voisinage ouvert $U$ de $X$ tel que $U \to Y$ soit
« normal » i.e.\ \emph{plat} à fibres géom.\ normales.

\emph{Dém} Lorsque on suppose $\mathcal{O}_{Y,y}$ \emph{excellent}
[il suffit : à fibres formelles géom.\ normales]

a) Itou les hyp.\ \add{restent} vraies au voisinage \add{de $x$}
\struck{\ill{}}.

\begin{tikzcd}[column sep=small, row sep=small]
X' \arrow[r, no head] \arrow[d] & X \arrow[d, no head] \\
Y' \arrow[u] \arrow[ur] \arrow[r, no head] & Y
\end{tikzcd}

\note{croquis, sous « EGA IV 9 » ; points $x_0$, $x'_0$ sur $X'$, $x$,
$x'$ sur $X$, $y$, $y'$ sur $Y$, $y_0$, $y'_0$ sous $Y'$ ; à côté un
second croquis semblable, où $X'$ et $Y'$ sont reliés à $X$, $Y$ par
des traits obliques, avec $Y''$ \uncertain{et} $y''_0$, $y_0$ ; les
directions des traits sans flèche ne sont pas marquées}

\emph{Vérifions} l'une des \ill{} : les pts de $X'$ où $X'$ est
réduit \ill{} \ill{} avec \ill{} chaque fibre.
\uncertain{Remarquons} que l'ens.\ \add{ouvert $V$} des pts où $X/Y$ est
\struck{plat} « séparable » i.e.\ \emph{plat} à fibres séparables est,
\struck{\ill{}} dans un voisinage de $x$, dense dans chaque fibre.

$V$ est un ouvert : EGA IV 12

$V$ contient les pts de $X_y$ où $X_y$ est réduit EGA IV 15.2.3
\marginal{donc $V \cap X_y$ est dense dans $X_y$ et on conclut par
équidimensionnalité}\note{au crayon, en bas à gauche}

\page{11}
\note{un long trait au crayon traverse la page en diagonale, de haut en
bas ; le texte reste lisible dessous}

\emph{NB} \struck{\ill{}} Supposons $S = \mathrm{Spec}\, k$,
\struck{\ill{}} $k$ corps alg.\ clos, et $X = S$,
\struck{\ill{} \ill{}} $S_0 =$ schémas \add{de t.f.}\ étales sur $X$
$=$ \uncertain{revêtements} \uncertain{locaux} sur $X$,
$\mathfrak{S} =$ schémas de t.f.\ sur $S$,
$\mathfrak{M} =$ imm.\ ouvertes.
\note{les trois égalités sont disposées l'une sous l'autre ; les
lettres lues $\mathfrak{S}$, $\mathfrak{M}$ sont des capitales de forme
gothique}

La catégorie $\mathcal{C}$ est celle des schémas \struck{\ill{}} finis
sur $k$ \add{\uncertain{\& connexes}}, avec la top.\ de Zariski, un
faisceau \add{dessus} revient à la donnée d'un \struck{\ill{}} foncteur
$\Lambda \mapsto F(\Lambda)$, $\Lambda$ fini sur $k$, transformant
\ill{} \ill{} en produits. Les faisceaux qu.\ coh.\ \add{sur
$\mathrm{Spec}\, \Lambda$} sont ceux \ill{} \ill{} \struck{\ill{}} de la
forme $\Lambda' \mapsto M \otimes_\Lambda \Lambda'$, avec $M$ un
$\Lambda$-module. \struck{\ill{}} \uncertain{En revanche}, les faisceaux
de prés.\ finie sont ceux \ill{} \ill{} \struck{\ill{}} $M$ \add{de
t.f.}. Mais \ill{} le \ill{} \ill{} pour lesquels \ill{} \ill{} \ill{}
\ill{} \ill{} \uncertain{soit} par exemple
\[
  \mathcal{O}_\Lambda^n \xrightarrow{\ u\ } \mathcal{O}_\Lambda
\]
défini par \struck{\ill{}} $u_1, \ldots, u_n \in \Lambda$, alors
$\mathrm{Ker}\, u$ est un « faisceau sur \struck{$\mathrm{Spec}$}
$\Lambda$ »\,?
\[
  \Lambda' \mapsto \mathrm{Ker}\bigl(u_{\Lambda'} \colon
  \Lambda'^{\,n} \to \Lambda'\bigr)
\]
qui n'est pas en général déduit de $\mathrm{Ker}\, u$ comme
$\Lambda' \mapsto (\mathrm{Ker}\, u) \otimes_\Lambda \Lambda'$,
manque de platitude \ldots

\note{au bas de la page, des croquis et notes au crayon sans rapport
visible avec ce qui précède et qui reprennent les pages 3--4 :
$T' \to T$ au-dessus de $X' \to X$ (points $x'$, $x$), $T'_y$ et
$T_y$ ; « $U \cap T_y$ dense dans $T_y$ », « $U \cap T'_y$ dense dans
$T'_y$ », « $T' = \bigcup T'_i$, $T = \bigcup T_i$ » ; et à l'encre, à
droite, des signes $\hat{I}$, $\hat{\eta}$,
$\hat{X} \cup \overline{\{x\}}$, dont l'articulation n'est pas lisible}

\page{12}
\note{au crayon ; suite de la démonstration de la page 10 (après le a)
de celle-ci)}

b) \emph{On suppose que les hyp.\ sont satisfaites en tous les pts de
$X$}.

$U$ l'ouvert où $f$ est \emph{séparable}, il est dense sur chaque fibre.
\add{$\overline{V}$ (adh.\ schématique)}

À prouver que $X = U$ et plat sur $Y$.

On avait \uncertain{remarqué}, on il \uncertain{revient} à prouver
$X_y$ est séparable sur $k(y)$, or soit $y'$ une gén.\ maximale de $y$,
et \uncertain{appliquons} un $Y' \to Y$, donc $X'_{y'_0}$ est réduit,
donc $X'$ \uncertain{plat} sur $Y'$, donc
$\overline{(\ )}$ \add{(adh.\ sch.)}\ $= X'_{y'_0}$ est aussi
réduit. Donc Hironaka s'applique, $X'_{y'_0}$ est normale, donc $X_y$
aussi\,: cqfd

c) Pour prouver la platitude de $X/Y$, on applique le \emph{critère
valuatif de platitude d'une adhérence schématique}.
\note{les indices $y'_0$ sont peu sûrs, et la phrase « donc
$\overline{(\ )}$ (adh.\ sch.)\ $= X'_{y'_0}$ » est écrite
serrée sous la précédente ; « critère valuatif \ldots\ schématique » est
souligné. Le reste du feuillet est blanc}

\page{13}
\note{un trait oblique au crayon traverse le haut de la page}

Considérons p.\ ex.\ \struck{\ill{}} le faisceau d'anneaux
\uncertain{standard} $\mathcal{O}$ sur $\mathcal{C}$, et considérons un
module $F$ \add{quasi-cohérent} sur $\mathcal{O}$.
\struck{En \ill{} $F_{\mathcal{E}}$,} \struck{\ill{}} Alors $F$ induit
sur chaque $\mathcal{M}_{/T}$ (pour $\overset{X}{\uparrow} T_0 \to T$)
un faisceau \emph{quasi-cohérent} $\overset{X}{\underset{T}{\uparrow}}$
avec \uncertain{isomorphismes} en l'objet \uncertain{variable}
$(T_0 \to T)$ de $\mathcal{C}$, et de façon \emph{cartésienne} \ldots\
\struck{De plus} \ldots
\marginal{à vérifier}\note{dans la marge gauche, avec un trait
vertical ; les petits $X$ surmontant les flèches sont peu sûrs}

En particulier on voit que l'on obtient un « faisceau formel » sur le
\ill{} \ill{} \ill{} complété formel d'un $X$ \struck{\ill{}} $Z$ le
long de \ill{} $Z$, \ill{} $\delta \colon X \to Z$, \ill{} $X$
[\uncertain{fini} sur un \ill{} \ill{} \ill{} $\in \mathrm{Ob}\,
\mathcal{C}$] et on \uncertain{a} les \uncertain{vérifications}
\uncertain{fin.}\ \ill{} \ill{} un \uncertain{vérification}
\uncertain{essentielle} \uncertain{obtiendra} \ill{}, par un
\uncertain{argument} facile
\[
  H^*(\mathcal{C}_{/\delta}, F) \simeq H^*(\widehat{Z}, F|\widehat{Z})
\]
où le \uncertain{deuxième} \ill{} est pris comme \ill{} \uncertain{limite}
\ill{} de la catégorie des systèmes projectifs \uncertain{adiques} sur le
complété formel $\widehat{Z}$ de $Z$ le long de $X$ \ldots
\marginal{\uncertain{précision} à \uncertain{vérifier}}\note{dans la
marge gauche, en face de la formule, avec un trait vertical}

\page{14}
\note{page de notes et de croquis au crayon, sans texte suivi ; elles
accompagnent le Th.~1.4 (pages 3--4). Les éléments en sont donnés dans
l'ordre de la page, sans chercher à les relier}

$F$, \quad $F_U \to G_U$, \quad $X \supset U$ ;
\quad $F \to i_*(G_U)$ se factorisant par $G$ ;
\quad $X' \supset U'$, $G'$, au-dessus de $Y' \to Y$ ;
\quad points $\tilde{x}_1$, $\tilde{x}$, $\tilde{x}'$.

\[
  A = \varprojlim \bigl(A_\alpha = A/\mathfrak{J}^{\alpha+1}\bigr)
\]
\[
  \left.
  \begin{array}{l}
    M \to N \\
    M_\alpha \xrightarrow{\ \sim\ } N_\alpha
  \end{array}
  \right\rbrace
  \Longrightarrow
  \forall\, \mathfrak{p} \in \mathrm{Spec}\, A,\
  \mathfrak{J} \not\subset \mathfrak{p},\
  M_{\mathfrak{p}} \xrightarrow{\ \sim\ } N_{\mathfrak{p}}
\]
\note{$f$ (modules de type fini, semble-t-il) est écrit en indice à
côté de $A$, de $M_\alpha$ et de $N_\alpha$ ; sous la conclusion, une
parenthèse peu lisible avec $\widehat{M}_{\mathfrak{p}}$}

$A \to B$, $\mathfrak{A} \subset A$, $\mathfrak{A}B$ ;
$(\mathfrak{A}B \cap A) \uncertain{=} \mathfrak{A}$ ;
$A/\mathfrak{A} \overset{?}{\hookrightarrow} B/\mathfrak{A}B$ ;
$A/\mathfrak{A}_\alpha \hookrightarrow B_\alpha$.

$x_\alpha \in \mathrm{Ass}\, G_U$ \add{($\Longrightarrow x_\alpha \in
\mathrm{Ass}\, G_y$)}, $T_\alpha = \overline{\{x_\alpha\}}$ ;
$x' \in (T_\alpha)_{y'}$ pt maximal
$\Longrightarrow x' \in \mathrm{Ass}\, G_{y'}$ (EGA IV 12.2) ; \emph{donc}
$x' \in U_{y'}$.
\note{la note « EGA IV 12.2 » est encerclée}

Contre-exemple \uncertain{(l'ouvert i-fini est $\simeq$)} :
$X \to Y$ birationnel, $U = Y$, $G = \mathcal{O}_X$, $T = X$,
$T_{y'}$ de dim.~0, $G' = \mathcal{O}_{U}$, $T_y$ de dim.~1.
\note{avec un croquis : une courbe $X$ au-dessus d'une droite $Y$, et
une fibre qui se relève en une composante verticale}

$x_\alpha \in \mathrm{Ass}\, G_U$ au-dessus de $y$,
$T_\alpha = \overline{\{x_\alpha\}}$ : $(T_\alpha)_{y'} \cap U$ est dense
dans $(T_\alpha)_{y'}$.

\page{15}
\note{au crayon ; la moitié inférieure est barrée de quatre longs traits
obliques}

Dém.\ de 1.5.

1.4.2 \quad \emph{Remarque 1} les conditions b), c) de 1.4
$\Longleftrightarrow$ pour \uncertain{tout} chgt de base $S' \to S$,
$S'$ \uncertain{artinien}, $\exists$ solution.

1.4.3 \quad \emph{Remarque 2} a) $\Longleftarrow$ $\exists$ solution après
tt chgt de base $S' \to S$, $S'$ trait.

\emph{Vrai théorème} \add{1.4} \quad \emph{Corollaire}
\[
  \exists \text{ solution} \Longleftrightarrow \exists \text{ solution
  après tt chgt de base } S' \to S,\ S' \text{ artinien ou trait}
\]
\marginal{\emph{Énoncé final}}\note{« Vrai théorème », souligné, est
écrit au-dessus de « Corollaire » et relié par un trait ; « Énoncé
final », souligné, est dans la marge gauche, avec un double trait
vertical}
\marginal{\uncertain{inutile} si les $\mathcal{O}_{Y,y}$ réduits}\note{encerclé,
relié à « trait »}

Pour prouver 1.5, il suffit donc de prouver \ill{} \ill{} \ill{}
$S' \to S$, $S'$ trait, $\exists$ solution. Donc \uncertain{on peut
supposer} $S$ est un trait.

\begin{tikzcd}[row sep=small]
X \supset U \arrow[d] \\
S
\end{tikzcd}

$F = G$, $G_U$. \quad $F = G$ est plat $/S$. Alors
$\forall x \in \mathrm{Ass}\, G_U$,
$x' \in \text{Maximal de } (T = \overline{\{x\}})_y \Longrightarrow
x' \in \mathrm{Ass}\, G_y$, donc \ldots

$x \in \mathrm{Ass}\, G_U$ au-dessus de $y'$, $y \in Y$,
$T = \overline{\{x\}}$, $T_y \neq \emptyset$ : les comp.\ de $T_y$ sont
de dim $\geqslant (\dim T_{y'}) = d$, donc ce sont des composantes de
$X_y$.

\page{16}
\note{page d'une autre main, qui n'est pas celle de Grothendieck et que
rien sur le feuillet n'identifie ; elle n'est pas recomposée. C'est une
mise au net en français de l'énoncé du Th.~1.4 : « Lemme 1.4' » --- les
notations et hypothèses étant celles de (1.1), $X$ et $Y$ localement
noethériens et $G = \mathrm{Im}\bigl(F \to i_*(G_U)\bigr)$ cohérent ; $Z$
l'espace sous-jacent à $X - U$ (une première formulation, « sous-schéma
réduit de $X$ ayant $X - U$ comme espace », est biffée), $y \in f(Z)$ ;
(d) pour tout $x \in Z_y = X_y \cap Z$, $\mathcal{O}_{X,x}$ à fibres
formelles géométriquement normales ; (a) pour tout
$x' \in \mathrm{Ass}(G_U)$, $T = \overline{\{x'\}}$, $T_y \cap U$ dense
dans $T_y$, et la même condition pour tout $X'$ étale sur $X$ ; (b) une
famille de morphismes locaux
$Y_i \to Y' = \mathrm{Spec}\, \widehat{\mathcal{O}}_{Y,y}$ après chacun
desquels le problème a une solution, l'intersection des noyaux de
$\widehat{\mathcal{O}}_{Y,y} \to \Gamma(Y_i, \mathcal{O}_{Y_i})$ étant
nulle ; (c) ou bien $f$ localement de présentation finie, ou bien les
$Y_i \to Y'$ finis. Conclusion : $G$ est plat sur $Y$ en tout point de
$X_y$ et $U_y$ contient $\mathrm{Ass}(G_y)$ ; en particulier, si les
conditions précédentes sont vraies en tout point $y$ de $f(Z)$, le
problème admet une solution cohérente. Aucune intervention de sa main
n'est sûre : les lettres $G$ de la conclusion (« Alors $G$ est plat sur
$Y$ ») sont repassées en gras, d'une encre plus noire, et un mot est
noirci sous le texte}

\page{17}
\note{suite de la page 13}

\begin{tikzcd}[column sep=small, row sep=small]
X \arrow[r, "\delta"] & Z \\
T_0 \arrow[u, "\varphi"] \arrow[r] & T \arrow[u, "u"'] \\
T_0^{(i)} \arrow[u] \arrow[r] & T^{(i)} \arrow[u, "u_i"']
\end{tikzcd}

\note{croquis de la marge gauche ; « $u_i$ » est suivi de
« $\in \mathcal{M}$ » encerclé}

\ldots\ explicitons les $H^q(\mathcal{C}_{/\delta_i}, F)$. Pour
simplifier la notation, on laisse tomber l'indice $i$.
\struck{Notons \ill{}} Supposons \struck{\ill{} \ill{}} que
$\delta \in \mathcal{N}$ [i.e.\ $Z \in \mathrm{Ob}\, \mathfrak{S}$].
\struck{donc} D'après un énoncé \ill{} \ill{} $\mathcal{C}_{/\delta}$ par
\uncertain{maintenant}, on \ill{} $(\mathcal{M}_{/\delta}$, i.e.\ \ill{}
catégorie des diagrammes

\begin{tikzcd}[column sep=small, row sep=small]
X \arrow[r] & Z \\
T_0 \arrow[u, no head] \arrow[r] & T \arrow[u]
\end{tikzcd}

qui sont des \uncertain{morphismes} \emph{cartésiens}
\add{$\mathcal{M}$-} dans $\mathcal{C}$, on \ill{} une catégorie
équivalente : $\mathcal{M}_{/Z}$, avec la topologie induite de
$\mathcal{C}$. \uncertain{On trouve}
\note{un trait vertical sépare le croquis de gauche du texte}

donc
\[
  H^q(\mathcal{C}_{/\delta_i}, F) \simeq
  H^q_{\mathcal{M},\mathcal{C}}(Z_n, F) \quad \ldots
\]
et \uncertain{finalement}
\[
  0 \to {\varprojlim_n}^{(1)} H^{q-1}_{(\mathcal{M},\mathcal{C})}(X^{(i)}_n, F)
  \to H^q(X^{(i)}, F) \to \varprojlim_n H^q_{(\mathcal{M},\mathcal{C})}(X^{(i)}_n, F)
  \to 0
\]
[sous l'hypothèse que les $X^{(i)}_n \in \mathrm{Ob}\, \mathfrak{S}$]
\note{« $Z_n$ » surcharge un mot biffé ; dans la suite exacte, encadrée,
l'exposant du premier terme se lit $q$ plutôt que $q-1$, et le terme de
droite surcharge une première écriture ; la flèche finale vers $0$ n'est
qu'amorcée. Deux traits verticaux traversent le bas de la page}

\page{18}
\note{quelques lignes en haut du feuillet, le reste est blanc}

\[
  \begin{array}{ccc}
    Z_U & \hookrightarrow & Z\,? \\
    \cap & & \cap \\
    U & \subset & X \\
    & \searrow & \downarrow \\
    & & Y
  \end{array}
\]

\note{croquis de la marge gauche, dessiné comme un triangle au-dessus de
$Y$ avec « plat » écrit à côté ; les inclusions $Z_U \subset U$ et
$Z \subset X$ sont des $\cap$ verticaux}

$Z_U = Z \cap U$ et tel que $Z_U$ soit universellement\,/$Y$ sch.\
dense dans $Z$.

\emph{NB} si $Z$ est \add{loc.\ de} prés.\ finie sur $Y$, ou $Y$, $Z$
noeth., \struck{l'hyp.\ \ill{}} la condition univoque
\[
  \Longleftrightarrow\ \forall y \in Y,\quad Z_y \cap U \supset
  \mathrm{Ass}\, Z_y \qquad \text{EGA IV 11.12\,?}
\]
\note{sous $Z_y \cap U$, un « $\cap$ » vertical et $Z_y$ ; « univoque »
est une lecture peu sûre}

\page{19}
\note{page de croquis et de formules, sans texte suivi ; elle reprend la
réduction de la page 4 (descente fpqc, hensélisé $X^h$, complété)}

\begin{tikzcd}[column sep=small, row sep=small]
X \arrow[d, no head] & X' \arrow[l, no head] \arrow[d, no head] & X'' \arrow[d, no head] & \\
Y & Y' \arrow[l] & Y'' \arrow[l] & Y'''  \arrow[l]
\end{tikzcd}

\note{au-dessus, $U$, $U'$, $U''$ et $F'$, $F'_{U'} \to G'_{U'}$,
$F''_{U''} \to G''_{U''}$ ; « fpqc » sous $Y' \to Y$ ; entre $Y'$ et
$Y''$ une double flèche, entre $Y''$ et $Y'''$ une triple, simplifiées
ici en flèches simples ; une ligne au crayon entoure la partie $X'$,
$X''$, $Y'$, $Y''$}

Au-dessous, un second croquis : $X' \to X$ fpqc, $F'$, $F$, au-dessus
de $Y$ ; puis $Z'_\alpha$, $Z'_i$ :
\[
  \left\lbrace
  \begin{array}{l}
    Z'_i \cap U'_y \neq \emptyset \\
    \text{i.e.\ } Z'_i \cap X'_y \neq \{x'\}
  \end{array}
  \right.
\]
et le carré

\[
  \begin{array}{llll}
    G_{U'} & U' \hookrightarrow X' = X^h & & X'_y \supset U'_y \\
    & \qquad\ \downarrow & & \ \downarrow \\
    G_U & U \subset X & & X_y \supset U_y
  \end{array}
\]

À droite :
\[
  \exists\, J = \{i_1, \ldots, i_n\} \subset I, \qquad
  \mathfrak{A}_n = \mathrm{Ker}\Bigl(\underbrace{\widehat{\mathcal{O}}_{Y,y}}_{A}
  \to \prod_{\alpha \in J} \mathcal{O}_{Y_\alpha, y_\alpha}\Bigr)
  \subset \widehat{\mathfrak{m}}_y^{\,n}
\]
OPS les $Y_i$ artiniens ;
\[
  A \to A/\mathfrak{A}_n \hookrightarrow \prod \mathcal{O}_{Y_\alpha, y_\alpha} .
\]
\note{« pour $n$ » est écrit au-dessus de $\exists$ ; la puissance
$\widehat{\mathfrak{m}}_y^{\,n}$ est encerclée ; une flèche courbe va de
la condition sur $Z'_i$ au produit. Au bas de la page, six croquis de
surfaces et de courbes dans un cube, deux d'entre eux marqués $Z'$,
$\Phi'$ et $\Phi$, $Z$, et un $Z_\alpha$ isolé ; ils ne sont pas
reproduits}

\page{20}
\note{un long trait oblique traverse la page de haut en bas}

Prenons une résolution injective $C^{\cdot}(F)$ de $F$, alors
$C^{\cdot}(F)_{S_\infty}$ est une résolution injective de $F_{S_\infty}$,
et les $C^{\cdot}(F)_{S_i} = C^{\cdot}(F)_{S_\infty} | S_i$ une
résolution injective de $F_{S_i} = F_{S_\infty} | S_i$. On a
\[
  \Gamma\bigl(S_\infty, C^{\cdot}(F)_{S_\infty}\bigr) =
  \varprojlim \Gamma\bigl(S_i, C^{\cdot}(F)_{S_i}\bigr)
\]
le système projectif des 2\textsuperscript{e} membres étant
\emph{strict} \struck{\ill{}} et \ill{} \ill{}.
\[
  H^n(S_\infty, F) \to \varprojlim H^n(S_i, F)
\]
et est bien sûr \emph{surjectif} ; il est bijectif si le système des
$H^{n-1}(S_i, F)$ vérifie ML. Dans le cas général, il semble qu'il faille
écrire
\[
  H^*(S_\infty, F) \Longleftarrow E_2^{pq} = {\varprojlim_i}^{(p)} H^q(S_i, F)
\]
mais avec ${\varprojlim}^{(p)} = 0$ si $p \neq 0, 1$, donc des suites
exactes
\[
  0 \to {\varprojlim_i}^{(1)} H^{n-1}(S_i, F) \to H^n(S_\infty, F)
  \to \varprojlim_i H^n(S_i, F) \to 0
\]
Dans le cas qui nous occupe, il vient :
\note{la page s'arrête sur ces mots ; la suite n'est pas dans ce lot.
L'exposant du premier terme de la suite exacte est surchargé}

\end{document}
